\section{Introduction}
\label{sec:introduction}

\begin{figure}[t]
\noindent\fbox{\parbox{0.97\linewidth}{\small
\textbf{Summary for student-support staff.}
\emph{The decision.} After the first semester, a support office can meet only about one in five
students and must decide whom to contact first.
\emph{What the model does.} It assigns each student a dropout-risk score from records the
university already holds (first-semester results, fees and background) and places the top 20\% on
a contact list.
\emph{Main findings.} On past data, the list had very high precision: 108 of the 109 flagged
test students dropped out. However, on validation data most flagged students had passed no units at all, which
staff can already see, and the list still misses about half of the students who later leave, partly
because only one in five can be contacted. A simpler, fully transparent model had comparable
final-test performance within the paired-bootstrap uncertainty. The ``reasons'' displayed are useful hints but may change
when recalculated, and there was no evidence that seeing them helped people predict the model's decision.
\emph{Limitations.} The model was tested only on Portuguese data from students with a known final
outcome (still-enrolled students were excluded). It does not explain \emph{why} a student is
struggling or what will help, and it must never be used to deny support.}}
\end{figure}

\subsection{Background and Problem Statement}

Student dropout is costly for students and institutions alike, and the dataset used here was created
specifically to help identify at-risk students early enough for support~\citep{realinho2021uci,
realinho2022predicting}. In practice, a
student-support office can hold only a limited number of outreach meetings after the first semester
and must therefore rank students and contact those at highest risk first. Ranking alone, however,
is insufficient: before a conversation, an adviser also needs to understand \emph{why} a student was
flagged, and those reasons must be trustworthy.

Tabular foundation models such as TabPFN~\citep{hollmann2025tabpfn} and
\tabicl~\citep{qu2025tabicl} have been reported to match or outperform tuned gradient boosting on
many small tabular benchmarks without dataset-specific training. They are, however, black boxes
whose explanations must be obtained post hoc, for example with SHAP~\citep{lundberg2017shap}, and
for high-stakes decisions an interpretable model should arguably be preferred whenever it performs
comparably~\citep{rudin2019stop}.

\textbf{Problem statement.} Benchmark accuracy does not establish that a foundation model improves a
capacity-limited outreach decision, nor that its post-hoc explanations are faithful, stable or
useful to the people who must act on them, and none of the representative studies in
Table~\ref{tab:prior-work} combines these three perspectives for dropout outreach. This project
therefore asks whether a foundation model's complexity is justified here and whether its
explanations withstand technical and human evaluation.

\subsection{Decision Setting}

The intended user is a student-support officer rather than the student. Predictions are made once,
after the first semester; second-semester information is excluded, as it would not yet be
available. The operating rule was fixed in advance: all students in a cohort are ranked by
predicted dropout probability, and exactly the top 20\% (the assumed outreach capacity) are flagged,
with ties broken by record order. A flag means ``offer a support conversation''; it never implies
withdrawing support or making an automated decision about the student. Errors carry different
costs: a missed student may leave without being offered help, whereas an unnecessary flag uses a
scarce meeting and may stigmatise a student who was coping, and a misread explanation may steer the
conversation towards the wrong problem.

\subsection{Research Questions}

\begin{enumerate}[nosep,leftmargin=*]
  \item[\textbf{RQ1.}] Does \tabicl add predictive value over simpler conventional and
  intrinsically interpretable baselines at a fixed 20\% outreach capacity?
  \item[\textbf{RQ2.}] Are its local explanations faithful, stable and appropriately constrained,
  and what do they reveal to an adviser about the students it flags?
  \item[\textbf{RQ3.}] Do explanation charts help people anticipate the model's outreach flag and
  identify useful support information?
\end{enumerate}

\subsection{Contributions}

\begin{enumerate}[nosep,leftmargin=*]
  \item A pre-specified, locked comparison of \tabicl with logistic regression, a depth-3 decision
  tree, XGBoost and a main-effects EBM at the same outreach capacity, with paired bootstrap
  intervals \textbf{(RQ1)}.
  \item A post-hoc audit of the explanations for all 544 validation students, covering
  faithfulness, stability and constrained recourse, together with an analysis of which students are
  flagged \textbf{(RQ2)}.
  \item A two-group forward-simulation study with 50 adult participants, including a qualitative
  analysis of 1,500 written answers \textbf{(RQ3)}.
  \item Sensitivity analyses for still-enrolled students, an admission-time feature window and a
  model without protected attributes, together with a local, validation-only research interface
  that illustrates how the outputs could be presented to staff.
\end{enumerate}

\subsection{Scope of the Work}

The study relies on one public dataset from a single Portuguese institution, a binary
Dropout-versus-Graduate target and one fixed train/validation/test split, and its fairness analysis
is descriptive. The human study involved volunteers acting as proxy users rather than support
officers; it therefore measures understanding of the model's output, not decisions about real
students or their outcomes. Table~\ref{tab:terms} defines the key terms.

\begin{table}[!htb]
\centering
\caption{Key terms used in this report.}
\label{tab:terms}
\footnotesize
\setlength{\tabcolsep}{3pt}
\begin{tabular}{@{}>{\raggedright\arraybackslash}p{2.0cm}p{6.0cm}@{}}
\toprule
\textbf{Term} & \textbf{Meaning} \\
\midrule
Outreach capacity & Share of students the office can meet: here the top 20\% of a cohort \\
Flag & The student is in that top 20\% and would be offered a meeting \\
Risk score & The model's estimated dropout probability; a score, not a certainty \\
Contribution (SHAP) & How far one group of information moved a student's score relative to a typical student; association, not cause \\
Faithfulness & Whether the reasons shown are those the model actually relies on \\
Stability & Whether the same top five reasons reappear when the explanation is recalculated \\
Recourse & A change to the record that would remove the flag; not advice \\
Forward simulation & People predict the model's decision from the same information \\
Locked final test & Held-out data used exactly once, after all choices were fixed \\
\bottomrule
\end{tabular}
\end{table}
